%% ============================================================
%% NOT PART OF THE HANDED-IN REPORT.
%% Working notes for structuring report.tex, with pointers to the
%% exact code locations you should read/cite while writing.
%% Delete this file before submission.
%% ============================================================


%% ================================================================
%% PART 0 — STATUS UPDATE / PASS 2 (2026-07-18)
%% ================================================================
%
% Since this notes file was first written, commit 4710124 ("Merge pull
% request #2 from bkarli/feature/new-gui", 2026-07-18) landed on main
% and REPLACED the frontend. This makes several things in Part 1/2
% below outdated. Summary of what changed and what it means for the
% report:
%
% 1. THE FRONTEND DESCRIBED BELOW NO LONGER EXISTS AS THE SHIPPED APP.
%    `blabber-gui` is still on disk but is no longer a workspace member
%    (see root Cargo.toml: members = ["blabber-core",
%    "blabber-app/src-tauri"]). The actual frontend is now
%    `blabber-app`, added whole in commit 5f1f89a ("add new GUI") and
%    iterated on through a889d47 ("add load spaces logic to join
%    space"). Every blabber-gui/... file path in Part 2 below is stale;
%    see the corrected pointers added inline.
%
% 2. IMPORTANT REGRESSION: voice calls and call rooms are backend-only
%    right now. blabber-app/src-tauri/src/voice_channel.rs exists but
%    its ENTIRE contents are commented out (start_call, hang_up,
%    my_endpoint_id, answer_call - all dead code), and there is no
%    call_room.rs at all in blabber-app. Compare to
%    blabber-gui/src-tauri/src/commands/voice_channel.rs and
%    commands/call_room.rs, where the same commands are live and wired
%    into invoke_handler!. blabber-app/src-tauri/src/lib.rs's
%    invoke_handler! only lists identity/space/room commands - no
%    voice, no call rooms. This has been true since the new GUI was
%    first added (5f1f89a) and is still true at HEAD.
%    => report.tex currently claims, as an ACHIEVED result, that users
%       can "start and end peer-to-peer voice calls" and describes a
%       whole group-call mesh architecture as part of the system. That
%       is only true of blabber-core (the Rust library) and of the old,
%       now-orphaned blabber-gui. If blabber-app is what gets handed
%       in, this is currently a factually wrong claim in "Supported
%       chat operations" and "Achieved results". Decide before
%       submission: (a) re-wire voice_channel.rs + add a call_room.rs
%       in blabber-app so the claim becomes true again, or (b) rephrase
%       the report to say voice/group-call is implemented at the
%       blabber-core level and demonstrated separately, but not yet
%       exposed in the new GUI. Don't leave the current wording as-is.
%
% 3. RESOLVED / NO LONGER APPLICABLE from Part 1's original gap list:
%    - "Persistence... not currently described in the report text at
%      all" is now WRONG as a note - report.tex already has a
%      "Persistence and application restart" subsection (added after
%      this notes file's first pass). No action needed there beyond
%      normal proofreading.
%
% 4. STILL OPEN / CONFIRMED UNCHANGED at HEAD (re-verified 2026-07-18):
%    - identity.rs:24 TODO ("mlock this one and use Zeroing") is still
%      present, secret is still a plain [u8; 32]. NOTE: root Cargo.toml
%      and blabber-core/Cargo.toml now declare a `zeroize = "1.9.0"`
%      dependency (workspace-level) that did not exist when Part 2 was
%      first written - but grepping the whole repo shows it is not
%      actually imported or called anywhere. Someone added the
%      dependency and never wired it up. Worth a one-line mention in
%      Limitations ("a zeroize dependency has been added but is not yet
%      used to clear the in-memory secret").
%    - node.rs perform_key_exchange_as_initiator/acceptor (now at
%      node.rs:460-467, diffie_hellman at 445-458) are still only
%      exercised by test_dh_key_exchange_produces_matching_keys; still
%      not consumed by VoiceChannel/MeshVoiceChannel to encrypt audio.
%      The caveat in Part 2's Security considerations entry still
%      holds verbatim.
%    - join_call_room/join_mesh mesh handshake (node.rs:387-442) is
%      unchanged: 16-byte room id + 1-byte ack, full mesh. Still
%      accurate.
%
% 5. Frontend architecture text in report.tex (Sec. "Core components" /
%    "Frontend architecture", describing App.vue switching between
%    login/chat/settings views via reactive state) describes the OLD
%    blabber-gui shape. Current blabber-app/src/App.vue is a two-line
%    shell (`<router-view />`) using vue-router; there are only two
%    top-level views (LoginView.vue, MainView.vue) with settings/room-
%    /space-creation implemented as dialogs and drawers
%    (IdentitySettingsDrawer.vue, CreateOrJoinSpaceDialog.vue,
%    CreateRoomDialog.vue) rather than a separate SettingsView. This
%    paragraph needs a rewrite, not just a proofread.


%% ================================================================
%% PART 1 — PROPOSED ACADEMIC STRUCTURE
%% ================================================================
%
% Your current report.tex already has good bones. Academic project
% reports (esp. for a DPI/distributed-systems course) are usually
% expected to follow roughly this arc:
%
%   1. Title page / Abstract
%   2. Introduction
%      1.1 Motivation / Background
%      1.2 Goals of the project
%   3. Problem statement / Requirements
%      2.1 Functional requirements
%      2.2 Non-functional / security requirements
%   4. Related work / Background technology
%      3.1 Why peer-to-peer
%      3.2 The Iroh protocol stack (endpoint, docs, gossip, blobs)
%      3.3 CRDTs (could also live later, see below)
%   5. System architecture
%      4.1 High-level overview (GUI / backend / networking layers)
%      4.2 Backend architecture (Node as entry point)
%      4.3 Core domain model (Identity, Space, Room, Invite, CallRoom)
%      4.4 Frontend architecture (Vue + Tauri commands)
%   6. Implementation
%      5.1 Identity creation & encrypted storage
%      5.2 Space & room creation, document model
%      5.3 Synchronization workflow (join, sync, live updates)
%      5.4 Conflict-free replication (CRDTs via Iroh Docs)
%      5.5 Voice communication (1:1 P2P calls)
%      5.6 Group voice calls (mesh architecture)
%   7. Security considerations
%      6.1 Threat model
%      6.2 Encrypted identities
%      6.3 Authentication via cryptographic identity
%      6.4 Message confidentiality/synchronization
%      6.5 Voice call key exchange (Diffie-Hellman)
%   8. Evaluation / Achieved results
%      7.1 Supported operations
%      7.2 Requirements coverage (map back to problem statement!)
%      7.3 Known limitations
%   9. Conclusion & future work
%  10. (Optional) References / Bibliography
%
% Notes on what to ADD vs. what you already have:
%
% - You already have sections for most of this (Introduction, Problem
%   statement, System architecture, Synchronization workflow, CRDTs,
%   Security considerations, Supported chat operations, Achieved
%   results, Conclusion). What's structurally missing for "academic
%   standard":
%     * An Abstract (short, 5-8 sentences, before Introduction).
%     * A "Related work / Background" section introducing Iroh's
%       building blocks (Endpoint, Docs, Gossip, Blobs) BEFORE you use
%       them in System architecture — right now Iroh Docs/Gossip/Blobs
%       are named in 3.1-3.2 without explanation of what each one does.
%     * A "Limitations / Future work" subsection — right now Conclusion
%       is a stub. Academic reports should be explicit about what does
%       NOT work / was not finished (e.g. TODO in identity.rs about
%       zeroizing secrets, missing message encryption at rest, no
%       group text-chat permissions model, no NAT traversal discussion).
%     * A explicit mapping from "Problem statement" requirements (a-f,
%       and the 4 security requirements) to what was actually
%       implemented — this is usually expected as a table in the
%       Evaluation section, and you already have almost all the
%       necessary facts in "Achieved results", just not tabulated
%       against the original requirement list.
%     * Bibliography/references for Iroh, Argon2, ChaCha20-Poly1305,
%       CRDT literature (Shapiro et al.) — currently zero citations
%       despite naming several protocols/algorithms.
%     * Figures: at least one architecture diagram (GUI <-> Tauri <->
%       Node <-> Iroh) and one sequence diagram for the synchronization
%       workflow (Section 4) would strengthen the report considerably;
%       you already have tikz loaded (see report.tex line 19-21) but
%       unused.


%% ================================================================
%% PART 2 — WHERE TO FIND THE DEFINITIONS FOR EACH CHAPTER
%% ================================================================
%
% Format: [Report section] -> [file:lines] -> what it gives you
%
% --- System architecture / Backend architecture (Node) ---
%   blabber-core/src/node.rs:30-45   struct Node fields (identity,
%       endpoint, gossip, router, blobs, docs, author, events,
%       spaces, active_call_rooms) -- this IS your "entry point"
%       paragraph, one sentence per field.
%   blabber-core/src/node.rs:346-359 Node::run() -- the exact boot
%       sequence: create_endpoint -> create_gossip -> create_blobs ->
%       create_docs_engine -> create_author -> create_router.
%       Use this as the backbone of an architecture diagram / the
%       "Node" paragraph (report.tex line 74).
%   blabber-core/src/node.rs:88-99   create_endpoint() -- explains
%       "same endpoint every login" claim in report.tex Authentication
%       subsection (line 146-148): SecretKey::from_bytes(identity.secret)
%       is deterministic.
%   blabber-core/src/node.rs:173-179 create_router() -- shows the four
%       ALPNs accepted (GOSSIP_ALPN, BLOBS_ALPN, DOCS_ALPN, VOICE_ALPN,
%       CALL_ROOM_ALPN) -- good concrete evidence for "networking layer"
%       paragraph (report.tex line 72).
%
% --- Core components / Identity ---
%   blabber-core/src/identity.rs:25-29  struct Identity fields.
%   blabber-core/src/identity.rs:45-51  derive_key() -- Argon2 usage,
%       cites report.tex line 140-142 directly.
%   blabber-core/src/identity.rs:54-93  store() -- ChaCha20-Poly1305
%       encryption, salt/nonce layout on disk (SALT_LEN=16, NONCE_LEN=12,
%       KEY_LEN=32 at lines 15-17) -- use these constants if you want
%       to describe the on-disk file format precisely.
%   blabber-core/src/identity.rs:97-139 load_from_disk() -- decryption
%       + manual byte-offset parsing of displayName/secret/author.
%       Worth a short "wire format" figure/table (byte offsets).
%   Note: line 24 has "TODO: mlock + Zeroing" -- good material for your
%       Limitations section (secret is a plain [u8;32], not zeroized).
%
% --- Core components / Space ---
%   blabber-core/src/space.rs:44-55   struct Space (id, name, info Doc,
%       members Doc, rooms Vec<Room>, docs handle) -- matches report.tex
%       line 88-89 almost verbatim; cite fields directly.
%   blabber-core/src/space.rs:60-89   Space::new() -- creates two Iroh
%       docs (info, members) and calls insert_self_as_member.
%   blabber-core/src/space.rs:92-114  insert_self_as_member() -- shows
%       the Member struct/key scheme "member/{author}" -- use as example
%       of Iroh-doc-as-key-value-store pattern.
%   blabber-core/src/space.rs:117-146 Space::from_invite() -- imports
%       info/member docs from DocTicket -- ties to Invite (below) and to
%       report.tex Synchronization workflow line 111-115.
%   blabber-core/src/space.rs:188-210 create_room() -- shows RoomRecord
%       {id, name, ticket} stored under key "room/{uuid}" in the info doc
%       -- this is the concrete mechanism behind "the backend reads the
%       available room records" (report.tex line 113-114).
%   blabber-core/src/space.rs:214-247 sync_rooms() -- the exact loop that
%       reads room/ prefixed keys, deserializes RoomRecord, imports each
%       room's DocTicket, and starts watch() handles -- this is THE
%       function to reference for "Synchronization workflow" section.
%
% --- Core components / Room + live sync ---
%   blabber-core/src/room.rs:20-26   struct Room (id, name, messages Doc,
%       cache).
%   blabber-core/src/room.rs:53-72   send_message() -- key scheme
%       "msg/{sent_at:020}-{author}" -- note the zero-padded timestamp
%       prefix, which is what gives lexicographic == chronological
%       ordering when listing via Query::single_latest_per_key(); good
%       detail for a "message ordering" paragraph.
%   blabber-core/src/room.rs:95-116  apply_event() -- reacts to
%       LiveEvent::InsertRemote/InsertLocal, updates in-memory cache,
%       and re-emits AppEvent::NewMessage -- this is the exact mechanism
%       behind report.tex line 117-121 ("Whenever an imported document
%       changes... sent as an application event").
%   blabber-core/src/room.rs:119-146 watch() -- wires Room::apply_event
%       into Node::watch_doc; a good target for a sequence diagram
%       (Iroh LiveEvent -> Room::apply_event -> broadcast::Sender
%       <AppEvent> -> Tauri event -> Vue listener).
%   blabber-core/src/node.rs:310-329 watch_doc() -- generic helper that
%       subscribes to any Doc and spawns a task calling a callback per
%       LiveEvent; explain this once, then say Room/CallRoom both reuse it.
%
% --- Core components / Invite ---
%   blabber-core/src/invite.rs:9-15   struct Invite (space_id, space_name,
%       info_ticket, member_ticket).
%   blabber-core/src/invite.rs:18-29  from_space() -- info doc shared
%       Read-only, members doc shared Write -- worth flagging explicitly:
%       this means every invited member can write directly into the
%       shared members document (permission-model detail relevant to
%       Security considerations).
%   blabber-core/src/invite.rs:31-47  serialize/deserialize -- postcard +
%       BASE32_NOPAD encoding -- explains the human-shareable "invite code"
%       string format mentioned in report.tex line 96-97 and 173.
%
% --- Synchronization workflow / persistence across restarts ---
%   blabber-core/src/node.rs:228-306 load_spaces() -- on-disk layout
%       root/<identity>/<uuid>/meta/invite.txt, re-imports each saved
%       invite and calls sync_rooms(); directly backs report.tex's
%       implicit claim that spaces persist across restarts (mentioned
%       in commit "Add persist and reload spaces across app restarts",
%       fcbcd2e) but NOT currently described in the report text at all
%       -- consider adding a short paragraph/subsection on persistence.
%   blabber-core/src/space.rs:157-164 create_directory() -- writes the
%       invite.txt used by load_spaces() above; the write side of the
%       same mechanism.
%
% --- CRDTs section ---
%   No CRDT algorithm is implemented directly in this repo (correctly
%   noted in report.tex line 130: "the application does not need to
%   implement its own... algorithms"). For citations/explanation, point
%   to Iroh's docs/replication model conceptually rather than to a file;
%   if you want a concrete "evidence" pointer, use:
%   blabber-core/Cargo.toml:17  iroh-docs dependency (fs-store feature)
%   -- and Query::single_latest_per_key() usage sites (space.rs:170,
%   217; room.rs:77; call_rooms.rs:68) which is the mechanism you get
%   "for free" from the CRDT doc (last-write-wins per key).
%
% --- Security considerations ---
%   blabber-core/src/identity.rs:45-93   Argon2 + ChaCha20-Poly1305 (see
%       above) -- Encrypted identities subsection.
%   blabber-core/src/node.rs:88-94        deterministic SecretKey ->
%       deterministic EndpointId -- Authentication subsection.
%   blabber-core/src/invite.rs (whole file) + space.rs:117-146 -- Message
%       synchronization subsection: only holders of a DocTicket can
%       import/sync, i.e. capability-based access control.
%   blabber-core/src/node.rs:429-452  diffie_hellman(),
%       perform_key_exchange_as_initiator/acceptor() -- Voice
%       communication subsection, cite directly: this is the actual DH
%       exchange over an Iroh bidirectional stream using x25519-dalek.
%   IMPORTANT CAVEAT to mention in Security considerations or
%       Limitations: grep shows the shared secret produced by
%       diffie_hellman() (node.rs:429-441) is computed but never
%       actually consumed by VoiceChannel/MeshVoiceChannel afterwards
%       (channel.rs sends raw f32 audio samples over
%       connection.send_datagram() with no encryption using that key,
%       see channel.rs:288-296 send_audio_chunk). Audio confidentiality
%       relies solely on Iroh's transport-level (QUIC/TLS) encryption,
%       not on an application-level cipher keyed by the DH secret. Also
%       corroborated by commit ed07da1 "removed de/encryption from
%       voicecall because this is already given via iroh". If your
%       report claims the DH exchange is *for* securing voice content,
%       double-check this is what you mean to write, or rephrase to say
%       transport security is provided by Iroh and DH is present but
%       currently unused/vestigial.
%
% --- Voice communication (1:1) ---
%   blabber-core/src/channel.rs:74-153   VoiceChannel (start_capture /
%       start_playback) -- CPAL device capture, mono-downmix, linear
%       resampling to a fixed WIRE_SAMPLE_RATE=48000 (line 16), then
%       chunked into datagrams sized to the connection's
%       max_datagram_size (samples_per_packet(), line 20-26).
%   blabber-core/src/channel.rs:155-196   ActiveVoiceCall -- capture/
%       playback lifecycle on a dedicated OS thread, Drop-based cleanup.
%   blabber-core/src/channel.rs:198-286   VoiceProtocol (ProtocolHandler
%       impl) -- accept() flow: ask GUI via on_incoming handler with a
%       30s timeout (line 252), then start ActiveVoiceCall, then wait on
%       either connection.closed() or an explicit hang-up Notify.
%   blabber-gui/src-tauri/src/commands/voice_channel.rs (whole file) --
%       Tauri command surface: start_call, hang_up, answer_call,
%       my_endpoint_id -- maps 1:1 to report.tex line 100-101 and to the
%       "Start and end peer-to-peer voice calls" bullet (line 177).
%
% --- Group voice calls (mesh architecture) — NOT YET IN REPORT TEXT ---
%   This is a whole feature (commits 11f26b4, 6ea19d6) with no
%   corresponding report section yet. You likely want a new subsection,
%   e.g. "5.6 Group voice calls" or extend "Voice communication".
%   blabber-core/src/call_rooms.rs:20-26   struct CallRoom (id, name,
%       participants, call_log Doc, cache) -- a call room is itself a
%       replicated-doc-backed entity, same pattern as Space/Room.
%   blabber-core/src/call_rooms.rs:54-79   log_call_started() /
%       list_call_log() -- call history persisted the same way chat
%       messages are (key prefix "call/").
%   blabber-core/src/node.rs:387-426   join_call_room() -- THE key
%       function to explain the mesh topology: connects directly to
%       every other known participant via CALL_ROOM_ALPN, exchanges a
%       16-byte room-id handshake (send room.id.as_bytes(), read 1-byte
%       ack), and adds accepted peers to a MeshVoiceChannel. This is a
%       full-mesh (every peer connects to every peer), NOT a relay/SFU
%       architecture -- an important architectural fact worth stating
%       explicitly (and a good "Limitations": O(n^2) connections, no
%       server-side mixing).
%   blabber-core/src/channel.rs:309-463   MeshVoiceChannel -- per-peer
%       jitter buffers (peer_buffers: HashMap<String, VecDeque<f32>>),
%       and start_playback()'s mixing loop (lines 405-432): every
%       MIX_CHUNK_MS=10ms (line 17) it drains one sample per active peer
%       buffer, sums, and divides by the active peer count -- this is
%       your concrete "audio mixing algorithm" paragraph, cite lines
%       405-432 directly and maybe include the averaging formula.
%   blabber-core/src/channel.rs:501-549   CallRoomProtocol (ProtocolHandler)
%       -- server-less signaling: on accept(), reads a 16-byte room id
%       from the incoming stream, looks it up in active_call_rooms, and
%       either accepts (adds peer to that room's MeshVoiceChannel) or
%       rejects with a 1-byte ack/nack.
%   blabber-gui/src-tauri/src/commands/call_room.rs (whole file) -- Tauri
%       surface: create_call_room, list_call_rooms, join_call_room,
%       leave_call_room.
%
% --- Frontend / GUI architecture ---
%   ** SUPERSEDED 2026-07-18: blabber-gui is no longer a workspace **
%   ** member (see PART 0). Use the blabber-app pointers below for **
%   ** anything you're writing NOW; the blabber-gui pointers are kept **
%   ** only for historical comparison (e.g. to describe the voice/call **
%   ** regression in PART 0 item 2). **
%
%   blabber-app/src-tauri/src/lib.rs:17-22   struct AppState -- now
%       just node, spaces, pending_call (no active_call,
%       incoming_call_handle, call_rooms, active_call_room fields --
%       those were dropped along with the commands that used them).
%   blabber-app/src-tauri/src/lib.rs:29-43   invoke_handler![...] list --
%       identity + space + room commands only (create_identity, login,
%       list_identities, logout, delete_identity, create_server,
%       list_servers, join_space, get_invite, create_room, list_rooms,
%       send_message, list_messages). No voice/call-room commands are
%       registered -- see PART 0 item 2.
%   blabber-app/src-tauri/src/event_bridge.rs -- spawn_event_bridge()
%       subscribes to Node's broadcast channel and re-emits every
%       AppEvent to the frontend via Tauri's `app_handle.emit(
%       "app-event", event)`; this is the concrete mechanism behind the
%       "sent through a Tokio broadcast channel...forwarded to the
%       frontend" claim in report.tex's Backend architecture paragraph.
%   blabber-app/src/App.vue -- just `<router-view />`; navigation is
%       now vue-router based, not manual view-switching.
%   blabber-app/src/views/{LoginView,MainView}.vue -- only two
%       top-level views now.
%   blabber-app/src/components/{layout,room,space,settings,chat}/*.vue
%       -- screen inventory for a UI walkthrough / screenshots section
%       (Achieved results, report.tex has a commented-out
%       "%Screenshots vom GUI" -- this is presumably where you intended
%       to add them). Notable components: SpaceSidebar.vue,
%       RoomSidebar.vue, CreateOrJoinSpaceDialog.vue,
%       CreateRoomDialog.vue, IdentitySettingsDrawer.vue,
%       MessageList.vue, MessageInput.vue.
%
%   --- historical (blabber-gui, no longer built) ---
%   blabber-gui/src-tauri/src/lib.rs:31-40   old AppState, included
%       active_call/call_rooms fields the new AppState dropped.
%   blabber-gui/src-tauri/src/commands/voice_channel.rs,
%   blabber-gui/src-tauri/src/commands/call_room.rs -- the LIVE
%       versions of the commands that are commented out / missing in
%       blabber-app; useful if you re-wire blabber-app before
%       submission (see PART 0 item 2, option a).
%   blabber-gui/src/API/tauriAPI.ts, blabber-gui/src/views/{LoginView,
%   ChatView,SettingsView}.vue -- old three-view frontend, replaced.
%
% --- Dependency / technology references (for a Bibliography) ---
%   Cargo.toml (workspace) lines 10-25 -- exact crate + version list:
%       iroh 1.0.2, iroh-blobs 0.103.0, iroh-docs 0.101.0,
%       iroh-gossip 0.101.0, argon2 0.5.3, chacha20poly1305 0.10.1 --
%       cite these exact versions if your course wants reproducibility
%       info.
%   blabber-core/Cargo.toml line 26  x25519-dalek 2 -- DH key exchange
%       library used in node.rs.


%% ================================================================
%% PART 3 — EXAMPLE CHAPTER DRAFTS (illustrative only, rewrite in
%% your own words / verify against the current code before use)
%% ================================================================

\iffalse   % wrap in \iffalse so this whole block is inert if you ever
           % accidentally \input this file from report.tex

\section{Related work and background technology}

Blabber is built on top of \textbf{Iroh}, a Rust networking stack for
building peer-to-peer applications without central infrastructure.
Iroh provides several independent building blocks that Blabber
composes:

\begin{itemize}
    \item \textbf{Iroh Endpoint} provides peer discovery and directly
    establishes QUIC connections between two peers identified by a
    public key (\texttt{EndpointId}), removing the need for a central
    connection broker.
    \item \textbf{Iroh Docs} provides replicated, multi-writer
    key-value documents. Every document is identified by a namespace
    and can be shared with other peers through a \texttt{DocTicket},
    which grants either read or write access.
    \item \textbf{Iroh Gossip} propagates document updates between
    peers that are interested in the same document, without requiring
    all peers to be directly connected to each other.
    \item \textbf{Iroh Blobs} stores the actual content referenced by
    document entries; a document entry only stores a content hash, and
    the associated bytes are fetched separately through the blob store
    (see \texttt{blobs().get\_bytes(entry.content\_hash())}, used
    throughout \texttt{space.rs}, \texttt{room.rs} and
    \texttt{call\_rooms.rs} whenever an entry is read back).
\end{itemize}

Blabber's backend (\texttt{blabber-core}) combines these four
components inside a single \texttt{Node} (see \texttt{node.rs}), and
adds one custom protocol of its own, \texttt{blabber/voice/0}, for
peer-to-peer audio transport, plus a second custom protocol
\texttt{blabber/callroom/0} for group-call signaling
(\texttt{channel.rs}).


\section{Persistence and application restart}

To make joined spaces available again after the application is
restarted, Blabber persists a small amount of local state per space:
the invite code used to join it. When a space is created or joined,
its invite is written to
\texttt{<identity>/<space-uuid>/meta/invite.txt}
(\texttt{Space::create\_directory}, \texttt{space.rs}). On startup,
\texttt{Node::load\_spaces} walks this directory, re-parses every
saved invite, re-imports the corresponding Iroh documents, and calls
\texttt{Space::sync\_rooms} to also restore every room that space
contains (\texttt{node.rs}). This means Blabber itself stores no chat
data or membership information outside of the replicated Iroh
documents; the only bespoke persisted artifact is the invite code
required to re-derive access to those documents.


\section{Group voice calls}

In addition to one-to-one calls, Blabber supports group voice calls
through a full-mesh topology: every participant establishes a direct
connection to every other known participant of the call room, rather
than routing audio through a relay or a selective forwarding unit
(SFU). A \texttt{CallRoom} (\texttt{call\_rooms.rs}) tracks the list of
known participant endpoint IDs and stores a call history in a
replicated Iroh document, following the same pattern used for chat
rooms.

When a user joins a call room (\texttt{Node::join\_call\_room},
\texttt{node.rs}), the local node connects to every participant it
already knows about via a dedicated ALPN
(\texttt{blabber/callroom/0}), sends the 16-byte call room identifier,
and waits for a one-byte acknowledgement before registering the
connection with a \texttt{MeshVoiceChannel}. Symmetrically,
\texttt{CallRoomProtocol::accept} handles inbound connections by
reading the same 16-byte identifier and looking up whether the local
node is currently active in that call room.

Because every peer receives one audio stream per other participant,
\texttt{MeshVoiceChannel} maintains one jitter buffer per peer
(\texttt{peer\_buffers}) and mixes them on a fixed 10ms interval by
averaging the buffered samples across all peers that currently have
data available (\texttt{MeshVoiceChannel::start\_playback},
\texttt{channel.rs}). This full-mesh design avoids any server-side
component but scales the number of connections quadratically with the
number of participants, which we discuss further in
Section~\ref{sec:limitations}.


\section{Limitations and future work}
\label{sec:limitations}

\begin{itemize}
    \item The \texttt{Identity} secret key is held in a plain
    \texttt{[u8; 32]} in memory and is not zeroized on drop (see the
    \texttt{TODO} comment in \texttt{identity.rs}); a memory-scraping
    attacker could recover it while the application is running.
    \item The Diffie-Hellman key exchange performed before a
    one-to-one voice call (\texttt{node.rs},
    \texttt{perform\_key\_exchange\_as\_initiator/acceptor}) derives a
    shared secret that is not currently used to encrypt the audio
    payload; confidentiality of the audio stream instead relies
    entirely on the transport-level encryption already provided by
    Iroh's underlying QUIC connection.
    \item Group voice calls use a full-mesh architecture
    (\texttt{Node::join\_call\_room}), which does not scale well beyond
    a small number of simultaneous participants, since every
    participant must maintain a direct connection and audio stream to
    every other participant.
    \item [... add any remaining gaps you are aware of, e.g. no
    read-receipts, no message editing/deletion, no rate limiting,
    no NAT traversal fallback discussion, etc.]
\end{itemize}

\fi


%% ================================================================
%% PART 4 — ACADEMIC-STANDARD LINE-LEVEL REVIEW (pass 2, 2026-07-18)
%% ================================================================
%
% Went through report.tex sentence by sentence against the current
% code. Below are concrete, fixable issues, grouped by type. Line
% numbers refer to report.tex as it currently stands.
%
% --- A. Typos / grammar (quick fixes, no judgement calls needed) ---
%   L46  "goverment" -> "government".
%   L46  "Government Id's" -> "government IDs" (no apostrophe for a
%        plural; "ID" not "Id").
%   L53  "Keep track of messages and and be assured..." -- duplicated
%        "and".
%   L227 "persistant" -> "persistent".
%   L287 "%Screenshots vom GUI" -- stray German comment left in a
%        comment; either fill in the screenshots here or remove the
%        comment. Fine either way since it's a LaTeX comment and won't
%        render, but tidy it up before submission.
%   L294 "%Mehr persönlich?" -- same: a private author note left as a
%        comment inside \section{Conclusion}. Resolve (write the
%        conclusion) and delete the comment.
%
% --- B. Register / parallelism ---
%   The requirements list (L50-59) mixes verb forms: "Create...",
%   "Keep track...", "Send...", "Receive...", but then "Being able to
%   start a call with another user" and "Navigate the user interface
%   while intuitively achieving their goals" break the imperative/
%   infinitive pattern set up by "Users should be able to:". Academic
%   style wants parallel structure across enumerated items. Proposal:
%   - (e) "Start a call with another user."
%   - (f) "Navigate the interface intuitively." (the "achieving their
%     goals" clause is redundant with "intuitively" and can be cut)
%
% --- C. Unsupported / unsourced factual claims ---
%   L46  "documents of at least seventy thousand users were leaked" --
%        a specific, checkable factual claim (presumably referring to
%        the 2023 Discord-related ID-verification vendor breach) with
%        no citation. Academic reports should either cite a source for
%        this number or soften it to something you don't need to
%        defend, e.g. "a large number of users' identification
%        documents were leaked" (still motivates the section, no
%        citation debt).
%   L46  "the Iroh protocol, which was introduced during the lecture"
%        -- fine as a factual/motivational aside, but if you want this
%        to read as an academic report rather than a diary entry,
%        consider citing the Iroh project (docs/website/GitHub) instead
%        of relying on "introduced during the lecture", which won't
%        mean anything to an external reader/grader outside your
%        specific course context.
%
% --- D. Structural / rigor gaps (expands on PART 1) ---
%   1. No Abstract. Every section currently jumps straight into prose
%      with no forward-referencing summary. Add a 5-8 sentence
%      abstract before Section 1.
%   2. Iroh's building blocks (Docs/Gossip/Blobs, Sec. 3.1) are
%      currently introduced INSIDE "System architecture", after they've
%      already been used unexplained in the System architecture intro
%      paragraph (L69-72: "Blabber uses Iroh Docs...Gossip...Blobs...")
%      -- the reader meets the terms before their definitions. Either
%      move 3.1 before the architecture overview paragraph, or move the
%      four building-block definitions into their own "Background"
%      section preceding "System architecture" (this matches the
%      Related-work-before-architecture convention noted in PART 1).
%   3. Zero citations in the whole document despite naming Argon2,
%      ChaCha20-Poly1305, Diffie-Hellman/x25519, CRDTs, and Iroh
%      itself. Minimum expected for a DPI report: one reference each
%      for Iroh, the CRDT concept (Shapiro et al. 2011), Argon2 (RFC
%      9106), and ChaCha20-Poly1305 (RFC 8439). Add a
%      \texttt{references.bib} + \texttt{\textbackslash bibliography}
%      or manual \texttt{thebibliography} block.
%   4. "Conclusion" (L293-296) is currently empty (just a
%      \texttt{\textbackslash newpage} and a leftover comment). This is
%      the single biggest missing piece for academic standard: a
%      report without a conclusion reads as unfinished regardless of
%      how good the preceding sections are. At minimum: restate what
%      was achieved, state explicitly which of the Section 2
%      requirements (a-f, plus the 4 security requirements) were met
%      vs. not (tie this to the requirements-mapping table proposed in
%      PART 1), and close with 2-3 sentences of future work pulling
%      from the Limitations list in PART 3.
%   5. No requirements-to-implementation traceability table. You have
%      6 functional + 4 security requirements in Section 2 and a prose
%      "Achieved results" in Section 8, but nothing that explicitly
%      maps requirement (a) -> implemented by X, satisfied Y/N. Add a
%      small table; this is usually the first thing a grader looks for
%      in a requirements-driven course project.
%   6. Given PART 0 item 2 (voice calls not wired in blabber-app), the
%      "Achieved results" and "Supported chat operations" sections
%      currently assert full voice support as fact. Whatever you decide
%      there (re-wire vs. rephrase), make sure Sections 7 and 8 agree
%      with each other and with what a grader can actually click
%      through in the delivered app -- an academic report is expected
%      to be verifiable against the artifact it describes.
%
% --- E. Minor style ---
%   The report uses "\\" line breaks to separate paragraphs throughout
%   (e.g. L45, L69-72, L105-106, ...) instead of a blank line (LaTeX's
%   native paragraph separator). This produces visually tight
%   "paragraphs" with no first-line indent removed / spacing adjusted,
%   which reads as a formatting artifact rather than an intentional
%   choice in a typeset report. Consider removing the \texttt{\\}s and
%   using blank lines (true \texttt{\textbackslash par}) instead, or if
%   you want the current tight spacing, wrap it in a consistent
%   \texttt{\textbackslash setlength{\textbackslash parskip}{...}}
%   rather than hand-inserting \texttt{\\} after every paragraph.
